\documentclass[10pt,a4paper]{amsart}
\usepackage[margin=19mm]{geometry}
\usepackage{amsmath,amssymb,mathtools}
\usepackage[scr=rsfs,cal=euler]{mathalfa}
\usepackage{xfrac}
\usepackage[unicode,hidelinks,bookmarks=false]{hyperref}
\newcommand{\ie}{i.e.\ }
\newcommand{\cf}{cf.\ }
\newcommand{\resp}{respectively\ }
\newcommand{\Subsection}{\subsection}
\providecommand{\nobreakdash}{\nobreak\hbox{-}}
\setlength{\emergencystretch}{2em}
\multlinegap0pt
\usepackage{bm}
\usepackage{bbm}
\usepackage{dsfont}
\usepackage{xspace}
\usepackage{cancel}
\usepackage{mathtools}
\usepackage{amsthm}
\makeatletter
\@ifundefined{theorem}{\newtheorem{theorem}{Theorem}}{}
\@ifundefined{lemma}{\newtheorem{lemma}[theorem]{Lemma}}{}
\makeatother
\newcommand{\func}[1]{\operatorname{#1}}
\title[Isospectral majorization and isoperimetric inequalities for ]{Isospectral majorization and isoperimetric inequalities for coherent states on the Bloch sphere}
\author{Luis Daniel Abreu}
\date{Source: arXiv:2608.12248v1 (CC BY 4.0)}
\begin{document}
\maketitle

\noindent\emph{Source and licence.} Isospectral majorization and isoperimetric inequalities for coherent states on the Bloch sphere. Version arXiv:2608.12248v1 (CC BY 4.0), distributed under the Creative Commons Attribution 4.0 International licence, as recorded in the arXiv licence metadata for this exact version and shown on the abstract page https://arxiv.org/abs/2608.12248v1. Full rights notice: licenses/arxiv-2608.12248-CC-BY-4.0.txt.\par\smallskip

\noindent\textit{Editorial scope.} Counted unit: the Lemma whose source label is \texttt{lem:classical} in the article (source0.tex lines 1146--1160, source label \texttt{lem:classical}), delivered together with the article's own complete original proof of it (source0.tex lines 1162--1241, one \texttt{proof} environment of 80 lines). No mathematics is written, repaired or re-derived: statement and proof are the article's own text line for line. Required context retained uncounted: eq:resolution (lines 121--124); lem:berezincalculus (lines 1067--1082); eq:betalimit (lines 1078--1081). Every definition, display and label of those lines is retained unchanged. Omitted from this package and not redistributed: 1--120, 125--1066, 1082--1145, 1161--1161, 1242--1608. Nothing inside a retained excerpt is omitted or altered: every retained body is delivered exactly as the article's lines are written, so no omission receipt of any kind accompanies this package. Citations: the retained excerpts carry no citation command at all, so no bibliography entry of the article is transcribed and no reference is redistributed. Reference adaptation: none was needed and none was made; every reference inside the delivered unit resolves inside it, so no unresolved reference or citation is produced. Printed slips kept literal: no printed word, symbol, hypothesis or number of the counted statement or of its proof was altered. Layout and class: the article's own class is \texttt{amsart} with packages including fontenc, lmodern, amsmath, amssymb, amsfonts, mathtools, microtype, hyperref, graphicx; the wrapper is the lane's pinned \texttt{amsart} editorial wrapper at 10pt on A4, which restates the theorem-like environments the retained text needs and adds the article's own macro definitions that the retained text needs (\texttt{\textbackslash func}), with extra wrapper packages (bm, bbm, dsfont, xspace, cancel, mathtools). The delivered numbering is the unit's own. Assets: no retained span contains a figure environment, a graphics-inclusion command, a PDF-page inclusion, a table or a TikZ picture; no image file is copied into this package, the package declares no asset dependency and no image file is redistributed with it. Licence: version arXiv:2608.12248v1 of this article is distributed under the Creative Commons Attribution 4.0 International licence, as recorded in the arXiv licence metadata for this exact version and shown on the abstract page https://arxiv.org/abs/2608.12248v1. A full rights notice is delivered at licenses/arxiv-2608.12248-CC-BY-4.0.txt. Characters: every retained excerpt is pure ASCII, so the delivered engine, XeLaTeX with its default Latin Modern text font, reports no missing character.
\begin{equation}
(N+1)\int_{\mathbb{C}}|\kappa _{N,z}\rangle \langle \kappa _{N,z}|\,dm(z)=I_{%
\mathcal{P}_{N}}\text{.}  \label{eq:resolution}
\end{equation}%

\begin{lemma}
\label{lem:berezincalculus} The lower symbol map $X\rightarrow Q_{X}$ is
injective and a linear isomorphism $End(\mathcal{P}_{M})$ onto $\mathcal{V}%
_{M}$. Moreover, the Berezin transform $B_{M}$ preserves each spherical
harmonic space and satisfies 
\begin{equation}
\mathcal{B}_{M}|_{\mathcal{Y}_{\ell }}=\beta _{M,\ell }I,
\label{eq:berezineigenvalue}
\end{equation}%
where $\beta _{M,\ell }>0$ for $0\leq \ell \leq M$. Thus $B_{M}^{-1}$\ is
defined on $\mathcal{V}_{M}$. Moreover, for fixed $\ell $, 
\begin{equation}
\beta _{M,\ell }\longrightarrow 1\quad \text{as }M\longrightarrow \infty .
\label{eq:betalimit}
\end{equation}
\end{lemma}

\begin{lemma}
\label{lem:classical} Let $X_{M}$ be operators on $\mathcal{P}_{M}$ such
that 
\begin{equation*}
0\leq X_{M}\leq I,\qquad Q_{X_{M}}=q,
\end{equation*}%
where $q$ is the lower symbol of an operator on a fixed space $\mathcal{P}%
_{N}$. Then, for every continuous function $\Phi :[0,1]\rightarrow \mathbb{R}
$, 
\begin{equation}
\lim_{M\rightarrow \infty }\frac{1}{M+1}\func{Tr}\Phi
(X_{M})=\lim_{M\rightarrow \infty }\int_{[0,1]}\Phi (t)d\mu _{M}(t)=\int_{%
\mathbb{C}}\Phi (q(z))dm(z).  \label{eq:classicalconv}
\end{equation}
\end{lemma}

\begin{proof}
By Lemma \ref{lem:berezincalculus}, the function $q$ belongs to $\mathcal{V}%
_{N}$. For $M\geq N$, define 
\begin{equation*}
a_{M}=\mathcal{B}_{M}^{-1}q
\end{equation*}%
on this fixed finite-dimensional harmonic space. Equations (\ref%
{eq:berezineigenvalue})--(\ref{eq:betalimit}) imply 
\begin{equation}
a_{M}\longrightarrow q\quad \text{uniformly on the sphere}.
\label{eq:upperconvergence}
\end{equation}%
The lower symbol of $\mathcal{T}_{M}(a_{M})$ is $\mathcal{B}_{M}a_{M}=q$.
The map taking an operator to its coherent-state lower symbol is injective
by Lemma \ref{lem:berezincalculus}. Hence 
\begin{equation}
X_{M}=\mathcal{T}_{M}(a_{M}).  \label{eq:upperrepresentation}
\end{equation}%
It follows that 
\begin{align}
\frac{1}{M+1}\func{Tr}X_{M}^{2}& =\int_{\mathbb{C}}a_{M}(z)Q_{X_{M}}(z)%
\,dm(z)  \notag \\
& =\int_{\mathbb{C}}a_{M}(z)q(z)\,dm(z)\longrightarrow \int_{\mathbb{C}%
}q(z)^{2}\,dm(z).  \label{eq:secondmoment}
\end{align}%
This second moment will be enough to control the convergence for every
continuous $\Phi $ . Consider the spectral decomposition of $X_{M}$, 
\begin{equation*}
X_{M}=\sum_{j=0}^{M}\lambda _{M,j}|\psi _{M,j}\rangle \langle \psi _{M,j}|%
\text{.}
\end{equation*}%
On $\mathbb{C}\times \{0,\ldots ,M\}$, define the probability measure 
\begin{equation*}
d\pi _{M}(z,j)=|\langle \kappa _{M,z},\psi _{M,j}\rangle |^{2}\,dm(z)\text{,}
\end{equation*}%
where counting measure is understood in the discrete variable $j$.\ The $z$%
-marginal is $dm$, while (\ref{eq:resolution}) shows that the $j$-marginal
assigns mass $1/(M+1)$ to each index. For the variables 
\begin{equation*}
\Lambda _{M}(z,j)=\lambda _{M,j},\qquad Y_{M}(z,j)=q(z)\text{,}
\end{equation*}%
we have 
\begin{equation*}
\mathbb{E}_{\pi _{M}}(\Lambda _{M}\mid z)=q(z)=Y_{M}.
\end{equation*}%
Observing that 
\begin{equation*}
\mathbb{E}_{\pi _{M}}(\Lambda _{M}Y_{M})=\mathbb{E}_{\pi _{M}}\left( Y_{M}%
\mathbb{E}_{\pi _{M}}(\Lambda _{M}\mid z)\right) =\mathbb{E}_{\pi
_{M}}(Y_{M}{}^{2})
\end{equation*}

since the variables are real valued, we obtain:%
\begin{eqnarray*}
\mathbb{E}_{\pi _{M}}|\Lambda _{M}-Y_{M}|^{2} &=&\mathbb{E}_{\pi _{M}}\left(
\Lambda _{M}-Y_{M}\right) ^{2}=\mathbb{E}_{\pi _{M}}\Lambda _{M}{}^{2}-2%
\mathbb{E}_{\pi _{M}}\Lambda _{M}Y_{M}+\mathbb{E}_{\pi _{M}}Y_{M}{}^{2} \\
&=&\mathbb{E}_{\pi _{M}}\Lambda _{M}{}^{2}-\mathbb{E}_{\pi _{M}}Y_{M}{}^{2}
\\
&=&\frac{1}{M+1}\func{Tr}X_{M}^{2}-\int_{\mathbb{C}}q^{2}\,dm\longrightarrow
0\text{,}
\end{eqnarray*}%
where we used the limit (\ref{eq:secondmoment}) in the last line. Both
random variables take values in $[0,1]$. If $\omega _{\Phi }$ is the modulus
of continuity of $\Phi $, then for every $\delta >0$, 
\begin{align*}
\mathbb{E}_{\pi _{M}}|\Phi (\Lambda _{M})-\Phi (Y_{M})|& \leq \omega _{\Phi
}(\delta )+2\Vert \Phi \Vert _{\infty }\pi _{M}\{|\Lambda _{M}-Y_{M}|>\delta
\} \\
& \leq \omega _{\Phi }(\delta )+\frac{2\Vert \Phi \Vert _{\infty }}{\delta
^{2}}\mathbb{E}_{\pi _{M}}|\Lambda _{M}-Y_{M}|^{2}.
\end{align*}%
Now let $M\rightarrow \infty $ and then $\delta \downarrow 0$ to yield 
\begin{equation*}
\mathbb{E}_{\pi _{M}}\Phi (\Lambda _{M})-\mathbb{E}_{\pi _{M}}\Phi
(Y_{M})\longrightarrow 0.
\end{equation*}%
The two expectations are precisely the two sides of (\ref{eq:classicalconv}%
). 
\end{proof}

\end{document}
